The earliest presentations of Arithmetic Expression Geometry began with a
simple picture: a nested arithmetic expression can be read as a sequence of
updates of one current value, and that sequence can be drawn as a path.  This
picture was developed for one comb orientation.  The opposite comb orientation
was visible in the syntax but was not given an equally explicit geometry.
Division shows that the omission is substantive.  Dividing the evolving value
by a constant remains affine, whereas dividing a constant by the evolving
value introduces inversion, a pole, and the projective point at infinity.

This paper develops the geometry that the picture requires.  It is the first
paper of the series, and it is the canonical source for three things that later
papers import: the space in which arithmetic operations act as motions, the
dual reading of those motions, and the commutativization against which their
order defect is measured.  The order of the chapters follows the order in which
the objects are needed.

\paragraph{Expansions of an expression.}
Section~\ref{sec:p0-expansion} fixes the syntax.  A binary expression whose
internal operations have a unique temporal order can be written in two extreme
ways: the active subtree occupies the first operand slot, or the second.  We
call these the left- and right-expanded combs, and we prove the two elementary
facts on which everything else depends: each has a unique internal evaluation
order, and the two grammars are exchanged by passing to opposite operations.
The chapter then isolates the distributive law, which is the one algebraic
identity that relates two *different* expansions of the same expression, and
shows what it does geometrically: distributivity moves additive charge through
multiplicative charge, and it therefore produces two expansions with the same
operator and different charges.  That observation is the seed of the tearing
discussed in Section~\ref{sec:p0-acs-tearing}.  Mechanical expansion of
expressions, in the tradition of Wu Wen-tsun's mathematics mechanization
\cite{Wu2000MathematicsMechanization}, uses this law as a means of eliminating
the process; here the expansions themselves are the object of study.

\paragraph{Arithmetic expression spaces and arithmetic motion.}
Section~\ref{sec:p0-aes} introduces the geometric object of the series.  An
arithmetic expression space (AES) carries an assignment function, whose value at
a point is the current arithmetic value, and a metric, which measures the cost
of moving.  The four operations act as *motions*: addition and subtraction move
the value by a constant, multiplication and division scale it, and the motions
that realize them transport the assignment in exactly that way.  The basic
model $\mathfrak E_0$ is the upper half-plane with the invariant metric of the
positive affine group; its additive and multiplicative motions are affine maps
of the plane, and they satisfy an elementary Baumslag--Solitar-type relation
which says that a multiplicative motion transports the additive ruler.  The
chapter also gives the continuous form of the motion, its rectification, and the
first model with a puncture, $\mathfrak E_1$, which is a punctured disc whose
metric extends across the puncture while its assignment does not.

\paragraph{The projective dual.}
Section~\ref{sec:p0-ripples} reads the same operator word in the opposite
direction.  Instead of pushing one input forward, pull an entire output
condition backward; for a fractional linear map the pullback of a standard
horocycle is a line or a circle tangent at the operator's pole, and the family
of such pullbacks is what we call a ripple pencil.  Taken as a whole --- pole,
pulled-back family, and the fixed points of the same matrix --- we call this
dual reading \emph{ripple geometry}.  The name is descriptive and provisional:
the constructions are proved, their interpretation as a geometry intrinsic to
the arithmetic word is not, and Section~\ref{subsec:p0-ripple-status} states
exactly what is and is not claimed.
Sections~\ref{sec:p0-reciprocals} and~\ref{sec:p0-projective-unification}
complete the elementary picture, treating reciprocals, poles, infinity, and
iterative fixed points, and then placing paths, poles, fixed points, and ripple
pencils in one matrix representation.

\paragraph{Commutativization, tearing, and contact.}
Section~\ref{sec:p0-affine-cocycles} computes the two translation coordinates of
an affine history, and Section~\ref{sec:p0-acs-tearing} constructs the
commutativization of the arithmetic motions: keep the total additive and
multiplicative charges of a history, forget the order in which they were
accumulated, and the result is a plane --- the accumulative commutative space,
ACS.  Two histories with the same charge bound a region of that plane, and the
exact value of their order defect is a weighted area of that region.  We read
this as a statement about the AES: the space tears by that amount, where
\emph{tearing} names a defect that is invisible in either motion separately and
becomes an area only in the commutativization (Definition~\ref{def:p0-tearing}).
Section~\ref{sec:p0-contact} gives the same defect a local form.  The three
coordinates (additive direction, multiplicative direction, current value) carry
a contact form whose kernel is a non-integrable horizontal distribution, and the
bracket of the two horizontal fields is the constant $\mu\lambda$ that also
appears as the density of the tearing.  This is the sense in which the contact
structure describes the tearing rather than merely accompanying it: it shows
that no local calibration can make both arithmetic directions straight at once.

\paragraph{Beyond linear language.}
Section~\ref{sec:p0-beyond-linear} is methodological, and its two propositions
are the technical reason for it.  The arithmetic tower generates an
infinite-dimensional Lie algebra, so no finite-dimensional group carries all its
motions; and every character of the commutativization agrees on
charge-compatible histories, so no character-based description can see tearing
at all.  We therefore use manifolds, metrics, and Lie groups where they are
proved to apply, and we do not adopt them as the ontology of the subject.  The
chapter records the alternative vocabulary that a companion programme, process
geometry, has begun to develop, together with the exact calibrations in which
that programme meets the laws proved here.  The proposals are labelled as
proposals.

\paragraph{Holed spaces and the language of tearing.}
Section~\ref{sec:p0-holed} closes the paper with the most recent development: a
definition of punctured arithmetic expression spaces, an exact computation of
the local shape of a puncture, and a register of six computed witnesses about
punctured configurations, each with its own status.  Their interpretive reading
--- that a hole is where a comparison must be performed --- is stated as a
motivation and not as a theorem, and the list of what the witnesses do not
establish is part of the chapter.

\paragraph{Status boundary.}
Everything proved in this paper is elementary algebra, upper-half-plane
geometry, one-dimensional dynamics, or a contact computation on a
three-dimensional local model.  Ordinary division remains a partial operation
even when its M\"obius continuation is well defined on the projective line.  No
convergence statement about infinite expressions is made beyond explicit
analytic or formal hypotheses.  The paper does not classify sequential trees,
does not develop function theory, and does not classify singular zero sets;
those belong to Paper~I, Paper~II, and Paper~III respectively, and
Section~\ref{sec:p0-interface} records the interfaces.  The migration that gives
this paper its present scope --- the geometric chapters of Paper~I now developed
here --- is recorded in the amendment
\texttt{governance/00b-paper-0-geometric-foundation-amendment.md}.
